% =====================================================================
% Adaptive Systems -- Assignment 2 report
% Homeostatic plasticity in CTRNNs: replicating Williams (2006) Ch.7
% and extending with a frozen-HP genetic-assimilation test.
%
% Max de la Nougerede, University of Sussex, MComp Adaptive Systems 2025/26.
%
% Compile: pdflatex -> biber -> pdflatex -> pdflatex
%
% WORD BUDGET (3000 words; methods, captions, bibliography, appendices
% are all EXCLUDED from the count per the spec):
%   Abstract        ~150
%   Introduction    ~700
%   Results/analyses ~1300
%   Discussion      ~800
%   (No conclusion -- spec says one is rarely needed.)
% Methods is unbudgeted: write it as long as reproducibility requires.
% =====================================================================

\documentclass[11pt,a4paper]{article}

\usepackage[margin=1in,top=0.9in,bottom=0.9in]{geometry}
\usepackage{amsmath, amssymb}
\usepackage{graphicx}
\usepackage{booktabs}            % professional tables
\usepackage{enumitem}
\usepackage[T1]{fontenc}
\usepackage{helvet}              % sans-serif for headings only
\usepackage[expansion=false]{microtype}
\usepackage{parskip}            % no indent, space between paragraphs
\usepackage{caption}
\usepackage[hidelinks]{hyperref}

% Sans-serif section headings (body stays roman).
\usepackage{sectsty}
\allsectionsfont{\sffamily}

% Bibliography: biblatex + biber, author-year science style.
\usepackage[english]{babel}
\usepackage{csquotes}            % recommended companion to biblatex
\usepackage[backend=biber, style=authoryear, sorting=nyt, maxcitenames=2, uniquelist=false]{biblatex}
\addbibresource{references.bib}

% Optional: uncomment if you typeset the GA pseudocode block in methods.
% \usepackage[ruled,vlined]{algorithm2e}

% Caption styling: bold label, small font.
\captionsetup{labelfont=bf, font=small, skip=6pt}

\title{\sffamily Homeostatic plasticity and the evolvability of CTRNN
ball-catching agents: a frozen-HP test for genetic assimilation}
\author{Max de la Nougerede}
\date{\today}

\begin{document}

\maketitle

% =====================================================================
% ABSTRACT  (~150 words)  -- write LAST
% =====================================================================
\begin{abstract}
% One short paragraph: what you did, why it matters, main result + conclusion.
% Lead with the assimilation verdict -- that is the finding that makes a
% selective reader want to read on.
\end{abstract}

% =====================================================================
% 1. INTRODUCTION  (~700 words)  -- write SECOND-TO-LAST
% =====================================================================
\section{Introduction}
\label{sec:intro}

% Planned thread (plant here, pay off in Discussion):
%  (a) Working definition of adaptivity -- Di Paolo grounded in Ashby.
%      State it explicitly; you must refer back to it in the discussion.
%  (b) HP background: biological (Turrigiano) + computational
%      substrate-level gains (Williams & Noble 2007).
%  (c) The substrate-vs-evolvability tension: HP helps the substrate, yet
%      Williams (2006) found online HP did NOT improve evolved fitness.
%  (d) Baldwin-effect frame: lifetime adaptation can guide evolution and
%      become genetically assimilated (Baldwin 1896; Hinton & Nowlan 1987;
%      Mayley 1996; Turney 1996). Dev-only = the Hinton & Nowlan setup;
%      behaviour-HP interleaves lifetime adaptation with fitness measurement.
%  (e) Stolting et al. (2023) as the candidate mechanism by which
%      assimilation can fail.
%  (f) Contribution statement: replication of Williams Ch.7 across four HP
%      conditions + four new analyses (3 of the controller, 1 of the search),
%      the central one being the frozen-HP / assimilation test.


% =====================================================================
% 2. METHODS  (UNBUDGETED -- excluded from word count)
% =====================================================================
% Section content lives in sections/methods.tex (kept separate so the
% word-counted body can be compiled/counted independently of methods).
% =====================================================================
% sections/methods.tex
% Methods section -- input by report.tex. Contains section content only:
% no preamble, no \documentclass, no \begin{document}.
%
% Spec constraints honoured here:
%  - reproducible at the level of equations + pseudocode; NO program code
%  - third-party algorithm (madvn/CTRNN) described as if our own
%  - every equation numbered and referred to by number in the text
%  - system diagram referenced (figure to be dropped in once drawn)
% Methods is EXCLUDED from the 3000-word count.
%
% [DIAGRAM] / [PSEUDOCODE] markers flag where a figure or algorithm block
% should be inserted. [CHECK] flags a value to confirm against source.
% =====================================================================

\section{Methods}
\label{sec:methods}

\subsection{System overview}
\label{sec:methods:overview}

The experimental system comprises four coupled components operating on
distinct timescales. A falling-shape environment presents circular
objects that descend through the agent's field of view. An agent body,
constrained to horizontal motion, carries an upward-facing fan of ray sensors
and attempts to position itself beneath each falling shape. A
continuous-time recurrent neural network (CTRNN) maps sensor readings to
motor commands. A homeostatic plasticity (HP) module adjusts the
network's weights and biases online, driving individual neurons towards a
target firing-rate range. Overarching these, a genetic algorithm (GA) evolves
the CTRNN's weights, biases and time constants across generations. The three
adaptive processes act on separate timescales: the network state changes every
timestep, HP modifies network parameters over the course of a trial, and the GA
modifies genotypes over generations. HP itself is not evolved; its parameters
are fixed across all runs. The experimental conditions vary when HP acts,
not whether it acts (Section~\ref{sec:methods:conditions}).

% [DIAGRAM] system_diagram.pdf -- the three coupled loops (sensorimotor,
% parameter-update/HP, selection/GA) on their three timescales. Spec asks for
% system diagrams explicitly; this is the place for it.
\begin{figure}[t]
  \centering
   \includegraphics[width=0.9\linewidth]{../figs/system_diagram.pdf}
  \caption{System diagram. The sensorimotor loop (agent body, ray sensors,
  CTRNN) runs every timestep; the homeostatic plasticity loop modifies network
  parameters within a trial; the genetic algorithm modifies genotypes across
  generations.}
  \label{fig:system}
\end{figure}

\subsection{CTRNN model}
\label{sec:methods:ctrnn}

The controller is a continuous-time recurrent neural network of $N = 5$
fully-connected nodes \parencite{beer1995}. Each node $i$ has an internal state
$y_i \in \mathbb{R}$ and a firing rate $z_i \in (0, 1)$. The state evolves
according to
%
\begin{equation}
  \tau_{y,i}\, \dot{y}_i \;=\; -y_i \;+\; \sum_{j=1}^{N} w_{ji}\, z_j \;+\; I_i,
  \label{eq:ctrnn-state}
\end{equation}
%
where $\tau_{y,i} > 0$ is the membrane time constant of node $i$, $w_{ji}$ is
the synaptic weight from presynaptic node $j$ to postsynaptic node $i$, $z_j$ is
the firing rate of node $j$ (Equation~\ref{eq:ctrnn-activation}), and $I_i$ is
the external input to node $i$, used to inject the sensor signal
(Section~\ref{sec:methods:body}). The firing rate is the logistic sigmoid of the
biased state,
%
\begin{equation}
  z_i \;=\; \sigma(y_i + b_i) \;=\; \frac{1}{1 + e^{-(y_i + b_i)}},
  \label{eq:ctrnn-activation}
\end{equation}
%
where $b_i \in \mathbb{R}$ is the bias of node $i$. Neuron gain is fixed at $1$
throughout, matching Williams's Chapter~7 specification, which does not evolve
gains \parencite{williams2006}.

Equation~\ref{eq:ctrnn-state} is integrated by the forward Euler method with a
fixed step size $\Delta t = 0.2$ \parencite{williams2006}:
%
\begin{equation}
  y_i(t + \Delta t) \;=\; y_i(t)
  \;+\; \frac{\Delta t}{\tau_{y,i}}
  \left( -y_i(t) + \sum_{j=1}^{N} w_{ji}\, z_j(t) + I_i(t) \right).
  \label{eq:ctrnn-euler}
\end{equation}

Of the five nodes, three are designated sensor nodes (receiving the three ray
sensors' signals as their external input) and two are motor nodes (whose firing
rates drive the body). There are no interneurons. This minimal architecture is
Williams's Chapter~7 specification; larger networks are explored in his later
experiments and are noted as future work (Section~\ref{sec:discussion}).
CTRNN integration used the \texttt{madvn/CTRNN} package \parencite{candadai2020},
whose algorithm is Equations~\ref{eq:ctrnn-state}--\ref{eq:ctrnn-euler}.

\paragraph{Genotype encoding.}
A genotype is a vector of $N^2 + 2N = 35$ real-valued alleles, each in
$[-1, 1]$, mapped linearly onto the phenotypic parameter ranges of
Table~\ref{tab:param-ranges}. The first $N^2 = 25$ alleles encode the weight
matrix, the next $N = 5$ encode the biases, and the final $N = 5$ encode the
time constants.

\begin{table}[t]
  \centering
  \caption{Phenotypic parameter ranges. Each genotype allele lies in $[-1,1]$
  and is mapped linearly onto the corresponding range.}
  \label{tab:param-ranges}
  \begin{tabular}{llc}
    \toprule
    Parameter & Symbol & Range \\
    \midrule
    Weights        & $w_{ji}$    & $[-10, 10]$ \\
    Biases         & $b_i$       & $[-10, 10]$ \\
    Time constants & $\tau_{y,i}$ & $[1, 4]$ \\
    \bottomrule
  \end{tabular}
\end{table}

\subsection{Homeostatic plasticity}
\label{sec:methods:hp}

Homeostatic plasticity adjusts each neuron's incoming weights and its bias so as
to drive the neuron's firing rate towards a target range $[H_L, H_U]$
\parencite{dipaolo2000, williams2006}. For each node $i$, a plastic facilitation
term $\rho_i \in [-1, 1]$ measures how far the firing rate lies outside the
range:
%
\begin{equation}
  \rho_i(z_i) \;=\;
  \begin{cases}
    (H_L - z_i)/H_L & \text{if } z_i < H_L, \\[2pt]
    0 & \text{if } H_L \le z_i \le H_U, \\[2pt]
    (H_U - z_i)/(1 - H_U) & \text{if } z_i > H_U.
  \end{cases}
  \label{eq:rho}
\end{equation}
%
When the neuron fires below the lower bound, $\rho_i > 0$ and the rules below act
to increase its input drive; when it fires above the upper bound, $\rho_i < 0$
and they act to decrease it. Each afferent weight to node $i$ is then updated by
%
\begin{equation}
  \tau_w\, \dot{w}_{ji} \;=\; \rho_i\, \lvert w_{ji} \rvert,
  \label{eq:hp-weight}
\end{equation}
%
and each bias by
%
\begin{equation}
  \tau_b\, \dot{b}_i \;=\; \rho_i.
  \label{eq:hp-bias}
\end{equation}
%
The absolute value $\lvert w_{ji} \rvert$ in Equation~\ref{eq:hp-weight} ensures
that the update scales with the current magnitude of the weight regardless of its
sign: when $\rho_i > 0$, both excitatory and inhibitory weights grow in magnitude,
and when $\rho_i < 0$ both shrink, so it is the strength of a connection that is
regulated rather than its direction. Equations
\ref{eq:hp-weight} and \ref{eq:hp-bias} are integrated with the same forward
Euler step $\Delta t = 0.2$ as the network state. Within a timestep the order of
operations is: inject sensor input, advance the network state by
Equation~\ref{eq:ctrnn-euler}, recompute each $\rho_i$ from the updated firing
rates, then apply Equations~\ref{eq:hp-bias} and~\ref{eq:hp-weight}.

The homeostatic range and timescales take Williams's Chapter~7 values
(Table~\ref{tab:hp-values}). Note that Williams uses $[0.25, 0.75]$ for the
substrate-level analyses of Chapter~6 and in \textcite{williamsnoble2007}, but
switches to $[0.2, 0.8]$ in Chapter~7; we use the Chapter~7 values consistently,
including in the substrate-level sanity check
(Section~\ref{sec:results:substrate}), to avoid mixing parameter sets across the
project.

\begin{table}[t]
  \centering
  \caption{Homeostatic plasticity parameters (Williams Chapter~7 values).}
  \label{tab:hp-values}
  \begin{tabular}{llc}
    \toprule
    Parameter & Symbol & Value \\
    \midrule
    Lower homeostatic bound & $H_L$    & $0.2$ \\
    Upper homeostatic bound & $H_U$    & $0.8$ \\
    Synaptic scaling timescale & $\tau_w$ & $40$ \\
    Intrinsic plasticity timescale & $\tau_b$ & $20$ \\
    \bottomrule
  \end{tabular}
\end{table}

\subsection{Agent body and sensors}
\label{sec:methods:body}

The agent is a circular body of radius $5$ constrained to move horizontally. Its
horizontal position $x$ is driven by the difference between the two motor
neurons' firing rates,
%
\begin{equation}
  \tau_x\, \dot{x} \;=\; z_R - z_L,
  \qquad \tau_x = 0.2,
  \label{eq:body}
\end{equation}
%
where $z_L$ and $z_R$ are the firing rates of the left and right motor nodes.
The agent senses its environment through three ray sensors arranged in an
upward-facing fan spanning $\pi/6$ radians. Each ray returns a reading that
increases as an intersected shape draws nearer:
%
\begin{equation}
  S \;=\; S_{\max}\,\frac{D_{\max} - D}{D_{\max}},
  \qquad S_{\max} = 5,\;\; D_{\max} = 100,
  \label{eq:sensor}
\end{equation}
%
where $D$ is the distance along the ray to the nearest intersection with a
shape, and $D_{\max}$ is the sensor's maximum range; a ray that intersects
nothing within $D_{\max}$ returns zero. Each of the three sensor readings is fed
as external input $I_i$ to one of the three sensor nodes
(Equation~\ref{eq:ctrnn-state}).

\subsection{Environment and trial structure}
\label{sec:methods:env}

Each fitness trial presents a sequence of falling circular shapes. A shape
enters above the agent at a randomised horizontal offset and falls with a
randomised horizontal and vertical velocity, descending until it passes the
agent's vertical level. % [CHECK] shape radius and initial-height values
The agent's task is to position its body so that the falling shape passes
through it. At the start of each trial the network state, weights and biases are
reset to the genotype-specified values, so that no adaptation carries across
trials.

In the conditions that include a developmental phase
(Section~\ref{sec:methods:conditions}), each trial is preceded by a
developmental period of $6000$ timesteps during which the network
receives zero external input ($I = 0$) and HP runs freely. Williams does not
state the input regime during this period; we follow the Chapter~6 substrate-level
precedent, in which HP is run on isolated networks, and use zero input. The
period length ($1200$ simulation seconds at $\Delta t = 0.2$) is far longer than
a trial, consistent with a ``settle the network's intrinsic properties''
reading rather than a ``preview the task'' one.

\subsection{Fitness function}
\label{sec:methods:fitness}

Fitness rewards an agent for reducing its horizontal distance to each shape over
the course of that shape's fall and for finishing close to it. For a single
shape, let $D_0$ be the initial horizontal agent--shape distance and $D_f$ the
final distance at the moment the shape passes the agent. The per-shape score
combines two components, averaged equally:
%
\begin{equation}
  c \;=\; \tfrac{1}{2}\bigl(\phi_1 + \phi_2\bigr),
  \label{eq:per-shape-score}
\end{equation}
%
where $\phi_1$ is a relative improvement term,
%
\begin{equation}
  \phi_1 \;=\;
  \begin{cases}
    0 & \text{if } D_0 = 0, \\
    \operatorname{clip}\!\bigl(1 - D_f/D_0,\; 0, 1\bigr) & \text{otherwise,}
  \end{cases}
  \label{eq:phi1}
\end{equation}
%
measuring how much the agent reduced its fractional distance to the shape; and
$\phi_2$ is an absolute proximity term,
%
\begin{equation}
  \phi_2 \;=\; \operatorname{clip}\!\bigl(1 - D_f/S_{\max},\; 0, 1\bigr),
  \label{eq:phi2}
\end{equation}
%
measuring how close the agent finished relative to the worst-case reachable
separation. The normalising constant $S_{\max}$ is the maximum separation an
agent could reach given the shape's velocities:
%
\begin{equation}
  S_{\max} \;=\; \bigl(1 + \lvert v_{x,\text{shape}} \rvert\bigr)
  \cdot \frac{100}{\lvert v_{y,\text{shape}} \rvert},
  \label{eq:smax}
\end{equation}
%
where $v_{x,\text{shape}}$ and $v_{y,\text{shape}}$ are the shape's horizontal
and vertical velocities. Both components are in $[0, 1]$, so the per-shape
score is in $[0, 1]$. The trial fitness is the mean per-shape score over the
$20$ shapes in the trial, and an individual's fitness is the mean over $10$
independently-seeded trials.

\subsection{Genetic algorithm}
\label{sec:methods:ga}

The GA evolves a population of genotypes
(Section~\ref{sec:methods:ctrnn}). Selection is by tournament of size $K = 3$:
to produce each offspring, three individuals are drawn without replacement and
the fittest is chosen as a parent. The five fittest individuals from the current
generation are copied unchanged into the next generation (elitism of five). There
is no crossover. Mutation is applied per allele with probability $p_m = 0.03$; a
mutated allele is perturbed by Gaussian noise $\mathcal{N}(0, \sigma_m^2)$ with
$\sigma_m = 0.1$, and values that fall outside $[-1, 1]$ are reflected back into
range.

This GA matches Williams's elitism of five and per-allele mutation rate of $3\%$,
but replaces his rank-based roulette-wheel selection with tournament selection.
The consequences of this difference for comparison with Williams are discussed in
Section~\ref{sec:discussion}.

% [PSEUDOCODE] A short algorithm block for one GA generation would satisfy the
% spec's "pseudocode or flowcharts" preference cleanly. Suggested:
%   for each generation:
%     evaluate fitness of every individual (Sec. fitness)
%     copy the 5 fittest individuals unchanged to the next generation (elitism)
%     while next generation not full:
%       parent <- argmax fitness of K=3 individuals drawn without replacement
%       child  <- copy(parent)
%       for each allele: with prob p_m, allele += N(0, sigma_m^2); reflect to [-1,1]
%       add child to next generation
% Add via the algorithm2e or algorithmicx package if you want it typeset.

\subsection{Experimental conditions}
\label{sec:methods:conditions}

The four conditions differ only in when HP is active:
%
\begin{description}[leftmargin=2em, itemsep=2pt]
  \item[No HP.] HP never acts. The genotype-specified parameters are used
    unchanged throughout each trial.
  \item[Dev only.] HP runs during the developmental period before each trial
    and is then frozen for the trial itself: the developed weights and biases
    are held fixed while behaviour is evaluated.
  \item[Behaviour only.] There is no developmental period; HP runs from the
    genotype-specified parameters throughout the trial.
  \item[Both.] HP runs during the developmental period and continues to run
    throughout the trial.
\end{description}
%
``Frozen'' HP means the plasticity update is disabled while the current weights
and biases are retained at their adapted values; it does not reset parameters to
the genotype. Williams terms this adiabatic elimination. The same four
conditions define both the evolutionary runs and the frozen-HP test
(Section~\ref{sec:methods:analysis}).

\subsection{Analysis methods}
\label{sec:methods:analysis}

Four analyses are applied to the evolved populations; three examine the evolved
controllers and one the search itself. All operate on data saved during the
evolutionary runs and re-evaluations, and all plots are regenerated from saved
data rather than live runs.

\paragraph{Behavioural trajectories.}
For a representative individual per condition, the agent and shape positions and
the five neurons' firing rates are recorded at every timestep across a small set
of shared-seed trials, allowing direct comparison of the neural activity that
produces the catching behaviour.

\paragraph{Viable-range diagnostics.}
For each neuron at sampled generations, two quantities are computed from a
replayed trial: the fraction of timesteps the neuron's firing rate lies within
$[H_L, H_U]$ (denoted $\mathrm{frac}_V$), and an entry--exit rate counting
transitions between the in-range and out-of-range states per $1000$ timesteps.
Each is computed both with HP active and with HP frozen, so that the
contribution of ongoing plasticity to viable-range operation can be isolated.

\paragraph{Frozen-HP test.}
This test asks whether the behaviour evolved under HP has been genetically
assimilated, i.e.\ whether it persists once HP is removed. Each evolved
individual is evaluated twice: once under its training condition and once with HP
disabled. The three HP conditions differ in what disabling means. For
\emph{Dev only}, the frozen evaluation runs the individual entirely without HP:
no developmental phase is applied and the individual is evaluated straight from
its genotype. For \emph{Both}, HP runs during the developmental period as normal
but is then disabled (adiabatic elimination: the plasticity update is turned off
while the adapted weights and biases are retained) for the trial. For
\emph{Behaviour only}, HP runs for an initial window of the first ten shapes and
is then disabled by adiabatic elimination, with fitness measured over the
remaining shapes. \emph{No HP} serves as a measurement-noise control: because HP
was never active, disabling it should produce no fitness change. To establish the
noise baseline, each individual is additionally re-evaluated $20$ times with
different shape-sequence seeds under its training condition; a frozen-HP fitness
drop is compared against this within-individual variability.

\paragraph{Search dynamics.}
To characterise the larger adaptive system, the population's best fitness, mean
fitness and fitness spread (the standard deviation of fitnesses across the
population) are tracked across generations for each condition. The fitness
spread serves as a diversity proxy. Per-generation genotype arrays are not
retained in the run history; genotype-level diversity measures would require
post-hoc analysis of the per-generation best-individual archives, which was not
performed.



% =====================================================================
% 3. RESULTS AND ANALYSES  (~1300 words)
% =====================================================================
% Spec (EA clause): focus analyses on the LARGER adaptive system -- EA,
% problem space, fitness function, population, and the landscapes they
% codetermine. Build everything toward the two main points:
%   MP1: genetic assimilation has NOT occurred (frozen-HP test).
%   MP2: the Dev-only paradox -- best-performing yet most HP-dependent.
\section{Results and analyses}
\label{sec:results}

\subsection{Substrate-level sanity check}
\label{sec:results:substrate}
% 0.861 -> 0.477 (HP on) -> 0.539 (HP off). Establishes HP works and that
% some regulation persists, foreshadowing the frozen-HP test.

\subsection{Replication: four-condition evolvability}
\label{sec:results:replication}
% Fitness curves (mean +/-1 SD, faint runs) + final-fitness box plot.
% KW H=17.71, p=0.0005; Dev only sig. > all others (Bonferroni). Online-HP
% NOT worse than No HP -- divergence from Williams; flag as a discussion hook.
% Include the search-REPLICATION read here (early progress, run consistency).

\subsection{Behavioural trajectories}
\label{sec:results:trajectories}
% 2-4 representative panels (likely Dev only vs Behaviour only). Blue/cream/red
% = below/in/above viable range. Motivates the two 8b metrics.

\subsection{Viable-range diagnostics across evolution}
\label{sec:results:viablerange}
% frac_V and entry-exit rate, training vs HP-off. Summary table.
% Headline: no HP condition holds neurons in range without HP.

\subsection{Frozen-HP test: has assimilation occurred?}  % <-- MP1 lands here
\label{sec:results:frozenhp}
% Drops in baseline-SD units; controls near zero; collapse counts
% (5/10 Dev, 3/10 Behaviour, 3/10 Both). Verdict: assimilation has not
% occurred in any HP condition.

\subsection{Search dynamics of the larger adaptive system}
\label{sec:results:searchdynamics}
% Population best/mean/spread per generation. "Differ in WHERE they converge,
% not whether/how fast." Serves the spec's EA clause directly. Consider
% whether this should frame the section rather than close it.


% =====================================================================
% 4. DISCUSSION  (~800 words)
% =====================================================================
% Do NOT list the spec's questions and answer them in turn (spec warns
% against this). Weave. Build to MP1 + MP2. Required coverage:
%   - further work
%   - how adaptive is the system, per the intro's definition
%   - connection to the wider adaptive-systems area
%   plus at least one of the hypothesis-support / new-hypothesis options.
\section{Discussion}
\label{sec:discussion}

% Suggested flow:
%  1. Assimilation verdict as the headline (MP1): tie frozen-HP drop to the
%     Baldwin frame; HP is constitutive, not merely helpful.
%  2. The Dev-only paradox (MP2): best yet most HP-dependent -- explain via
%     the selection regime (GA never saw un-developed phenotypes).
%  3. Stolting mechanism as candidate explanation for assimilation failure.
%  4. Search-dynamics framing: HP changes the genotype-phenotype map across a
%     FIXED landscape (NOT a "moving landscape" -- Chris's correction).
%  5. Adaptivity revisited: which timescales are active in each condition;
%     HP as cybernetic negative feedback / Ashby ultrastability.
%  6. Divergence from Williams (online HP not worse) -- GA-conditioning artefact?
%  7. Limitations: single task, fixed tau_w/tau_b, n=10, no discrimination,
%     developmental-settling confound (one honest sentence).
%  8. Future work: discrimination; timescale-separation sweep; alternative
%     rho shapes; exact-Williams-GA comparison; larger networks.


% =====================================================================
% BIBLIOGRAPHY  (excluded from word count)
% =====================================================================
\printbibliography
\label{sec:refs}


% =====================================================================
% APPENDICES  (excluded; must NOT carry main-narrative material)
% =====================================================================
% \appendix
% \section{Supplementary figures}
% e.g. full per-neuron viable-range grids, full frozen-HP per-individual table.

\end{document}
